\documentclass{article}

\usepackage{microtype}
\usepackage{graphicx}
\usepackage{subcaption}
\usepackage{booktabs} 
\usepackage{hyperref}
\usepackage{amsmath}
\usepackage{amssymb}
\usepackage{mathtools}
\usepackage{amsthm}
\usepackage[capitalize,noabbrev]{cleveref}
\usepackage{algorithm}
\usepackage{algorithmic}

% Use the following line for the initial blind version submitted for review:
\usepackage{icml2026}

\renewcommand{\theHalgorithm}{\arabic{algorithm}}

\icmltitlerunning{Sharpness-Aware Test-Time Adaptive Merging for Robust Multi-Task Learning}

\begin{document}

\twocolumn[
  \icmltitle{Sharpness-Aware Test-Time Adaptive Merging for \\ Robust Multi-Task Learning}

  \begin{icmlauthorlist}
    \icmlauthor{Anonymous Author}{equal,yyy}
  \end{icmlauthorlist}

  \icmlaffiliation{yyy}{Department of Computer Science, University of Machine Learning, Location, Country}
  \icmlcorrespondingauthor{Anonymous Author}{anon@ml.edu}

  \icmlkeywords{Model Merging, Test-Time Adaptation, Sharpness-Aware Minimization, Robustness}

  \vskip 0.3in
]

\printAffiliationsAndNotice{} 

\begin{abstract}
  Model merging is an emerging, training-free paradigm that combines multiple task-specific expert models into a single multi-task network. Recent work has advanced this by employing Test-Time Adaptation (TTA) to dynamically adjust merging coefficients and classifier heads on unlabeled test batches, achieving impressive specialized performance. However, adapting model parameters on extremely small, uncurated test batches introduces a severe vulnerability: the optimization is highly prone to overfitting, culminating in sharp, brittle minima that catastrophically fail under real-world distribution shifts or corruptions. To address this fundamental limitation, we propose \textbf{SATT-Merge} (\textbf{S}harpness-\textbf{A}ware \textbf{T}est-\textbf{T}ime Adaptive \textbf{Merge}), a novel framework that integrates Sharpness-Aware Minimization (SAM) directly into the test-time adaptation loop. By perturbing the joint space of merging coefficients and task-specific classification heads with worst-case adversarial noise, SATT-Merge guides the optimization toward flatter regions of the joint loss landscape. We introduce a tensorwise perturbation normalization to address the significant scale discrepancy between low-dimensional coefficients and high-dimensional weights. Empirically, SATT-Merge consistently outperforms existing state-of-the-art model merging baselines across multiple benchmarks under clean, noisy, and rotated conditions, demonstrating substantially improved robustness and generalization capabilities without requiring any extra supervision.
\end{abstract}

\section{Introduction}
\label{sec:intro}

Deep neural networks have achieved remarkable success across diverse domains, but training unified multi-task models remains a massive computational bottleneck. The standard practice of multi-task learning or sequential fine-tuning frequently suffers from catastrophic forgetting, gradient interference, or prohibitive computational costs. To circumvent expensive multi-task training, the machine learning community has turned to \emph{model merging} as a highly efficient and training-free alternative. Model merging techniques, such as Task Arithmetic \cite{yang2023adamerging, gargiulo2024task, he2024localizeandstitch}, linear mode connectivity, and Model Soups \cite{lotfy2023detecting, tran2025leveraging}, directly combine the weights of pre-trained expert models that share a common initialization. By doing so, they integrate diverse capabilities into a single consolidated model in a training-free manner.

While initial merging methods used uniform or static weights, recent advances introduce Test-Time Adaptation (TTA) to dynamically adjust merging coefficients on unlabeled data from the target domain. For instance, AdaMerging \cite{yang2023adamerging} optimizes coefficients by minimizing the entropy of model predictions on-the-fly. More recently, SyMerge \cite{symerge2024} proposes jointly adapting both the merging coefficients and task-specific classification heads using a self-labeled distillation loss, where the individual expert models serve as teacher guides to resolve task interference.

Despite their success, existing test-time adaptive merging methods suffer from a critical and under-explored flaw: \emph{overfitting during adaptation}. Because TTA relies on small, localized batches of unlabeled data (e.g., $16$ to $64$ samples) and operates in an unsupervised or self-labeled manner, the loss landscape is highly under-constrained. Optimizing merging coefficients and classifiers on these tiny slices of data typically converges to extremely sharp and narrow local minima. Under benign, in-distribution (ID) test sets, these sharp minima can yield high accuracy. However, in real-world scenarios, test data is often corrupted by noise, spatial transformations, or sensor degradation. Under these distribution shifts, the sharp minima found during test-time adaptation generalize extremely poorly, leading to severe representation collapse and catastrophic degradation of multi-task performance.

To bridge this gap, we draw inspiration from Sharpness-Aware Minimization (SAM) \cite{saim2026, lotfy2023detecting} and propose a flatness-seeking test-time adaptation strategy. In this work, we present \textbf{SATT-Merge} (\textbf{S}harpness-\textbf{A}ware \textbf{T}est-\textbf{T}ime Adaptive \textbf{Merge}), which integrates flatness regularization directly into the TTA loop of model merging. Rather than optimizing the self-labeled cross-entropy loss solely at the nominal weight parameters, SATT-Merge formulates a minimax objective that evaluates the loss over a neighborhood in the joint parameter space of merging coefficients and task-specific heads. Specifically, at each test-time step, we compute a worst-case adversarial perturbation, apply it to the parameters, evaluate the self-labeled soft cross-entropy loss, and perform the parameter update using gradients computed at this perturbed point. This explicitly regularizes the loss landscape's curvature on-the-fly, driving the merged model toward wide, flat minima.

Importantly, optimizing both merging coefficients (typically a $3$-dimensional vector in a $3$-task setup) and classification heads (containing tens of thousands of parameters) introduces a severe scale discrepancy. A global gradient normalization would cause the classification head gradients to dominate the perturbation, rendering the sharpness-aware objective ineffective for the merging coefficients. To resolve this, SATT-Merge employs a \textbf{tensorwise perturbation strategy} that normalizes the gradient and scales the adversarial step independently for each parameter tensor. This ensures that the flatness-seeking objective is applied equitably across both the macro-merging level (coefficients) and the micro-representational level (classifier weights).

Our contributions are summarized as follows:
\begin{itemize}
    \item We identify and analyze the overfitting vulnerability of test-time adaptive model merging under out-of-distribution (OOD) corruptions, showing that standard self-labeling adaptation produces brittle, sharp minima.
    \item We propose \textbf{SATT-Merge}, the first framework to integrate Sharpness-Aware Minimization (SAM) with test-time adaptive model merging. SATT-Merge optimizes both merging coefficients and classification heads on a minimax objective to seek flatter, more robust loss regions.
    \item We introduce a \textbf{tensorwise perturbation normalization} strategy to address the severe scale and dimension discrepancies between merging coefficients and neural weights, enabling stable and effective joint TTA.
    \item We conduct extensive experiments across sequential and independent TTA protocols under clean, Gaussian noise, and spatial rotation corruptions. SATT-Merge consistently outperforms state-of-the-art baselines (including Task Arithmetic, AdaMerging, and SyMerge), achieving substantial performance improvements under severe distribution shifts.
    \item We present a comprehensive hyperparameter sensitivity analysis of the perturbation radius $\rho$, empirically demonstrating a clear ``Goldilocks'' effect and confirming the tight link between test-time flatness and OOD robustness.
\end{itemize}

\section{Related Work}
\label{sec:related}

\textbf{Model Merging.} Model merging has emerged as a key area of study for combining multiple specialized neural networks without retraining. The foundational work in this area includes Model Soups \cite{lotfy2023detecting}, which averages weights of models trained on different hyperparameter configurations, and Task Arithmetic \cite{yang2023adamerging, gargiulo2024task, he2024localizeandstitch}, which adds task-specific "task vectors" to a base pre-trained model. Other notable works include RegMean \cite{wang2024localizing}, Fisher Merging, and ZipIt, which align representations before merging. To balance performance across diverse tasks, modern model merging focuses on finding optimal coefficients, which are often solved using validation datasets or optimization heuristics. Recently, OrthoMerge \cite{orthomerge2026} proposed merging models on the non-Euclidean Riemannian manifold of the orthogonal group using Lie algebra representations to decouple weight updates into rotational and residual components. While OrthoMerge focuses on preserving representation geometry through architectural constraints, SATT-Merge focuses on regularizing the optimization landscape itself during unsupervised test-time adaptation.

\textbf{Test-Time Adaptation (TTA).} TTA adapts a pre-trained model to a target distribution during inference using only unlabeled test data. Early works like TENT \cite{lee2024entropy} minimize Shannon entropy on batch predictions, while subsequent methods like CoTTA address error accumulation. In the context of model merging, AdaMerging \cite{yang2023adamerging} adapts merging coefficients at test-time using entropy minimization. SyMerge \cite{symerge2024} advances this by jointly updating the merging coefficients and classification heads via a self-labeled distillation setup, where specialized expert models provide soft pseudo-labels. However, these methods are highly susceptible to overfitting on small test batches, especially under corruption, causing representation drift.

\textbf{Sharpness-Aware Minimization.} Generalization in deep learning is closely tied to the flatness of the loss landscape at the converged minimum. Sharpness-Aware Minimization (SAM) \cite{saim2026} formalizes this by finding parameters that lie in a region of uniformly low loss, achieved by adding a gradient-based perturbation during training. Recent works have applied flatness-seeking objectives to continual learning (e.g., SAIM \cite{saim2026}), which uses a sharpness-aware block coordinate descent (SA-BCD) fine-tuning method to locate flatter minima. SATT-Merge differs from SAIM fundamentally: while SAIM optimizes flatness during the supervised fine-tuning phase of expert models, SATT-Merge operates entirely on-the-fly during the unsupervised test-time adaptation phase of merged models, addressing the unique overfitting vulnerabilities of small-batch TTA.

\section{Background \& Problem Formulation}
\label{sec:background}

Consider a pre-trained base model $f_{\theta_0}$ parameterized by $\theta_0 \in \mathbb{R}^D$, representing a shared initialization. We have a set of $T$ expert models $\{f_{\theta_i}\}_{i=1}^T$, where each expert has been fine-tuned from $\theta_0$ on a specific task $i$. The parameter update for each expert is defined as the task vector $\tau_i = \theta_i - \theta_0$.

In multi-task model merging, the goal is to construct a single merged encoder $f_{\theta(\lambda)}$ whose parameter vector $\theta(\lambda)$ is defined as:
\begin{equation}
    \theta(\lambda) = \theta_0 + \sum_{i=1}^T \lambda_i \tau_i
\end{equation}
where $\lambda = [\lambda_1, \dots, \lambda_T]^\top \in [0, 1]^T$ are the merging coefficients. Each task $i$ is associated with a task-specific classification head $h_{\phi_i}$ parameterized by $\phi_i$.

During test-time adaptation, we are given a sequence of unlabeled test batches $\{X_i\}$ for each task $i$. In a self-labeled joint TTA setting (such as SyMerge), the individual expert models serve as frozen teachers to guide the adaptation of the merged model. Specifically, for a batch $X_i$ from task $i$, the frozen expert encoder $f_{\theta_i}$ and its head $h_{\phi_i^{\text{expert}}}$ produce soft teacher predictions:
\begin{equation}
    \hat{y} = \text{Softmax}\left( \frac{h_{\phi_i^{\text{expert}}}(f_{\theta_i}(X_i))}{\tau} \right)
\end{equation}
where $\tau$ is a temperature hyperparameter. The student merged model, consisting of the merged encoder $f_{\theta(\lambda)}$ and the task-specific student head $h_{\phi_i}$, outputs student predictions:
\begin{equation}
    \tilde{y} = \text{Softmax}\left( h_{\phi_i}(f_{\theta(\lambda)}(X_i)) \right)
\end{equation}
The joint parameters to be optimized are represented as $\mathbf{w} = [\lambda, \phi_1, \dots, \phi_T]$. The TTA objective is to minimize the soft cross-entropy loss between student and teacher predictions:
\begin{equation}
    \mathcal{L}(\mathbf{w}) = \frac{1}{T}\sum_{i=1}^T \mathcal{H}_{\text{soft}}(\tilde{y}, \hat{y})
\end{equation}
where $\mathcal{H}_{\text{soft}}$ represents the Kullback-Leibler divergence or soft cross-entropy. Optimizing $\mathcal{L}(\mathbf{w})$ directly on a small batch of test inputs $X_i$ leads to sharp, brittle minima that overfit the local batch features, severely degrading performance when the inputs are subjected to out-of-distribution corruptions.

\section{Proposed Method: SATT-Merge}
\label{sec:method}

We introduce \textbf{SATT-Merge}, which addresses the overfitting issue of test-time adaptive model merging by incorporating a sharpness-aware objective. SATT-Merge guides the test-time adaptation of the joint parameters $\mathbf{w}$ toward flatter regions of the self-labeled loss landscape, improving robustness against distribution shifts and corruptions.

\subsection{Sharpness-Aware Test-Time Optimization}

Instead of minimizing the empirical loss $\mathcal{L}(\mathbf{w})$ at the single point $\mathbf{w}$, we define a minimax optimization problem over a neighborhood around $\mathbf{w}$:
\begin{equation}
    \min_{\mathbf{w}} \max_{\|\epsilon\|_2 \le \rho} \mathcal{L}(\mathbf{w} + \epsilon)
\end{equation}
where $\rho > 0$ is the perturbation radius that defines the size of the neighborhood. This objective minimizes the worst-case loss in the $\rho$-ball around the current parameter state, which prevents the optimization from landing in sharp, narrow valleys where a small perturbation in the parameters would cause a massive spike in loss.

To solve this minimax problem efficiently on-the-fly, we employ a first-order approximation. For each test batch $X_i$, we perform the following steps:

\textbf{1. Compute the worst-case perturbation $\epsilon$:}
Using a Taylor expansion of $\mathcal{L}(\mathbf{w} + \epsilon)$ around $\mathbf{w}$, the inner maximization can be approximated as:
\begin{equation}
    \epsilon^* \approx \arg\max_{\|\epsilon\|_2 \le \rho} \epsilon^\top \nabla_{\mathbf{w}} \mathcal{L}(\mathbf{w})
\end{equation}
The closed-form solution for this constrained optimization under the $\ell_2$-norm is:
\begin{equation}
    \epsilon^*(\mathbf{w}) = \rho \frac{\nabla_{\mathbf{w}} \mathcal{L}(\mathbf{w})}{\|\nabla_{\mathbf{w}} \mathcal{L}(\mathbf{w})\|_2}
\end{equation}
In practice, if the gradient norm is extremely small ($\le 10^{-12}$), we set $\epsilon^* = 0$ to prevent numerical instability.

\textbf{2. Perturb the parameters:}
We apply the computed perturbation to the joint parameter state:
\begin{equation}
    \mathbf{w}_{\text{pert}} = \mathbf{w} + \epsilon^*(\mathbf{w})
\end{equation}
This corresponds to perturbing both the merging coefficients $\lambda$ and the active student head parameters $\phi_i$.

\textbf{3. Compute gradients at the perturbed state:}
We evaluate the self-labeled soft cross-entropy loss at this worst-case perturbed point and compute the gradients with respect to the original parameters:
\begin{equation}
    g_{\text{pert}} = \nabla_{\mathbf{w}} \mathcal{L}(\mathbf{w}_{\text{pert}})
\end{equation}

\textbf{4. Restore parameters and update:}
Finally, we restore the original parameters $\mathbf{w}$ from $\mathbf{w}_{\text{pert}}$ and apply an optimization step (e.g., via Adam) using the perturbed gradient $g_{\text{pert}}$:
\begin{equation}
    \mathbf{w} \leftarrow \mathbf{w} - \eta \cdot \text{Adam}(g_{\text{pert}})
\end{equation}
where $\eta$ is the learning rate. We then clamp the merging coefficients $\lambda$ to the valid range $[0, 1]^T$ to ensure physical consistency.

\subsection{Tensorwise Perturbation Normalization}
A major challenge in joint test-time optimization of model merging is the vast discrepancy in dimensions and numerical scales between different parameter groups. Specifically, the merging coefficients $\lambda \in [0,1]^T$ represent macro-level blending weights of very small dimensionality, while the classification heads $\phi_i$ are deep neural network weights with high dimensions. If a global $\ell_2$ norm normalization is applied to the joint parameter gradient $g$, the high-dimensional weights $\phi_i$ completely dominate the norm, resulting in negligible perturbations applied to the merging coefficients $\lambda$.

To address this scale discrepancy, SATT-Merge employs a \textbf{tensorwise perturbation strategy}. We treat each parameter group (and each weight tensor within the classification heads) as a distinct optimization variable. For each independent parameter tensor $\mathbf{p} \in \mathbf{w}$, we compute its individual gradient $g_{\mathbf{p}} = \nabla_{\mathbf{p}} \mathcal{L}(\mathbf{w})$, and apply a tensorwise normalized perturbation:
\begin{equation}
    \epsilon_{\mathbf{p}}^* = \rho \frac{g_{\mathbf{p}}}{\|g_{\mathbf{p}}\|_2}
\end{equation}
By scaling and normalizing the perturbation independently per tensor, we ensure that both the macro-coefficients $\lambda$ and micro-weights $\phi_i$ are perturbed by an amount proportional to their respective gradient directions, preserving the sharpness-aware regularization effect across all scales of the multi-task model.

The complete algorithm is described in Algorithm~\ref{alg:satt}.

\begin{algorithm}[ht]
\caption{Sharpness-Aware Test-Time Adaptive Merging (SATT-Merge)}
\label{alg:satt}
\begin{algorithmic}[1]
\REQUIRE Base pre-trained state $\theta_0$, expert states $\{\theta_i\}_{i=1}^T$, frozen expert heads $\{h_{\phi_i^{\text{expert}}}\}_{i=1}^T$, test loaders $\{\mathcal{D}_i\}_{i=1}^T$, perturbation radius $\rho$, learning rate $\eta$, temperature $\tau$.
\STATE \textbf{Initialize:} Merging coefficients $\lambda \leftarrow [0.3, \dots, 0.3]^\top$, student heads $h_{\phi_i} \leftarrow \text{copy}(h_{\phi_i^{\text{expert}}})$.
\STATE \textbf{Define parameters to optimize:} $\mathbf{w} = [\lambda, \phi_1, \dots, \phi_T]$.
\FOR{task $i = 1, \dots, T$}
    \FOR{batch $X_i \in \mathcal{D}_i$}
        \STATE \textbf{Generate Teacher Soft Labels:} \\
        $\hat{y} \leftarrow \text{Softmax}\left( h_{\phi_i^{\text{expert}}}(f_{\theta_i}(X_i)) / \tau \right)$
        \STATE \textbf{Step 1: Compute Nominal Loss \& Gradients} \\
        $\theta(\lambda) \leftarrow \theta_0 + \sum_{j=1}^T \lambda_j (\theta_j - \theta_0)$ \\
        $\tilde{y} \leftarrow h_{\phi_i}\left(f_{\theta(\lambda)}(X_i)\right)$ \\
        $\mathcal{L} \leftarrow \mathcal{H}_{\text{soft}}(\tilde{y}, \hat{y})$ \\
        Compute nominal gradients: $g_{\mathbf{p}} \leftarrow \nabla_{\mathbf{p}} \mathcal{L}$ for all tensors $\mathbf{p} \in \mathbf{w}$
        \STATE \textbf{Step 2: Apply Tensorwise Perturbations} \\
        Initialize perturbations dictionary: $\mathcal{P} \leftarrow \emptyset$
        \FOR{tensor $\mathbf{p} \in \mathbf{w}$}
            \IF{$g_{\mathbf{p}}$ exists and $\|g_{\mathbf{p}}\|_2 > 10^{-12}$}
                \STATE $\epsilon_{\mathbf{p}} \leftarrow \rho \frac{g_{\mathbf{p}}}{\|g_{\mathbf{p}}\|_2}$
                \STATE $\mathcal{P}[\mathbf{p}] \leftarrow \epsilon_{\mathbf{p}}$
                \STATE $\mathbf{p} \leftarrow \mathbf{p} + \epsilon_{\mathbf{p}}$
            \ENDIF
        \ENDFOR
        \STATE \textbf{Step 3: Evaluate Loss at Perturbed Point} \\
        $\theta(\lambda) \leftarrow \theta_0 + \sum_{j=1}^T \lambda_j (\theta_j - \theta_0)$ \\
        $\tilde{y}_{\text{pert}} \leftarrow h_{\phi_i}\left(f_{\theta(\lambda)}(X_i)\right)$ \\
        $\mathcal{L}_{\text{pert}} \leftarrow \mathcal{H}_{\text{soft}}(\tilde{y}_{\text{pert}}, \hat{y})$ \\
        Compute perturbed gradients: $g_{\mathbf{p},\text{pert}} \leftarrow \nabla_{\mathbf{p}} \mathcal{L}_{\text{pert}}$
        \STATE \textbf{Step 4: Restore Parameters} \\
        \FOR{tensor $\mathbf{p} \in \mathbf{w}$}
            \IF{$\mathbf{p} \in \mathcal{P}$}
                \STATE $\mathbf{p} \leftarrow \mathbf{p} - \mathcal{P}[\mathbf{p}]$
            \ENDIF
        \ENDFOR
        \STATE \textbf{Step 5: Apply Parameter Update} \\
        $\mathbf{w} \leftarrow \mathbf{w} - \text{Adam}(g_{\text{pert}}, \eta)$
        \STATE \textbf{Constraint Check:} \\
        $\lambda_j \leftarrow \max(0, \min(1, \lambda_j))$ for all $j \in \{1,\dots,T\}$
    \ENDFOR
\ENDFOR
\end{algorithmic}
\end{algorithm}

\subsection{Theoretical Analysis and Scale Robustness}
\label{subsec:theory}

In this section, we provide a mathematical justification for why seeking flatness during test-time adaptation is essential for generalizing under out-of-distribution shifts, and formally analyze why a standard global SAM formulation fails in joint model-merging TTA.

\subsubsection{Generalization Guarantees on Unlabeled Target Domains}

During test-time adaptation, we optimize parameters $\mathbf{w}$ using a small, local batch $S = \{X_i\}_{i=1}^N$ drawn from the target domain $D_T$. Let $R(\mathbf{w}) = \mathbb{E}_{(x, y) \sim D_T} [ \ell(f_{\mathbf{w}}(x), y) ]$ represent the true target-domain multi-task risk, and let $\mathcal{L}_{\text{TTA}}(\mathbf{w}; S)$ represent our self-labeled soft cross-entropy surrogate loss on the batch $S$. 

By adapting PAC-Bayesian generalization bounds under parameter perturbations \cite{saim2026}, we can bound the true target domain risk with high probability.
\newtheorem{theorem}{Theorem}
\begin{theorem}
For any target distribution $D_T$, with probability at least $1 - \delta$ over the choice of the test-time batch $S \sim D_T^N$, the expected target-domain risk $R(\mathbf{w})$ satisfies:
\begin{equation}
\begin{aligned}
R(\mathbf{w}) \le & \max_{\|\epsilon\|_2 \le \rho} \mathcal{L}_{\text{TTA}}(\mathbf{w} + \epsilon; S) \\
& + c_1 \sqrt{\frac{\|\mathbf{w}\|_2^2 / \rho^2 + \log(N / \delta)}{N}} + \xi_{\text{dist}}
\end{aligned}
\label{eq:pac_bound}
\end{equation}
where $c_1 > 0$ is a complexity constant, and $\xi_{\text{dist}}$ is the representation and labeling discrepancy between the optimal target-domain classifier and the frozen expert teacher models.
\end{theorem}

Theorem 1 provides a direct theoretical foundation for SATT-Merge:
\begin{itemize}
    \item \textbf{Curvature Regularization:} The first term $\max_{\|\epsilon\|_2 \le \rho} \mathcal{L}_{\text{TTA}}(\mathbf{w} + \epsilon; S)$ is precisely the SATT-Merge inner optimization objective. Minimizing this term forces the optimization to find a parameter region where the loss is uniformly low, which mathematically minimizes the upper bound of the target-domain risk.
    \item \textbf{Small Batch Size Defense:} The complexity term scales with $O(1 / \sqrt{N})$. Since test-time adaptation operates on extremely small batches (e.g., $N=64$), the generalization gap can be exceptionally large. Under standard empirical risk minimization (e.g., SyMerge), the optimization is highly prone to finding a sharp local minimum that fits the local batch but generalizes poorly. SATT-Merge directly mitigates this risk by explicitly regularizing the curvature on-the-fly.
\end{itemize}

\subsubsection{Failure of Global Perturbations under Scale Discrepancies}

Joint test-time adaptation of model-merging optimizes a composite parameter vector $\mathbf{w} = [\lambda, \phi_1, \dots, \phi_T]$, containing both low-dimensional merging coefficients $\lambda \in [0, 1]^T$ (representing macro-level combination weights) and high-dimensional classification head weights $\phi \in \mathbb{R}^d$ (where $d \gg T$).

Suppose we apply a standard global $\ell_2$-bounded perturbation $\max_{\|\epsilon\|_2 \le \rho} \mathcal{L}_{\text{TTA}}(\mathbf{w} + \epsilon)$. The closed-form adversarial step is:
\begin{equation}
\epsilon^* = [\epsilon^*_\lambda, \epsilon^*_\phi] = \rho \frac{[g_\lambda, g_\phi]}{\|[g_\lambda, g_\phi]\|_2}
\end{equation}
where $g_\lambda = \nabla_\lambda \mathcal{L}_{\text{TTA}}$ and $g_\phi = \nabla_\phi \mathcal{L}_{\text{TTA}}$ are the gradients with respect to the merging coefficients and classifier heads, respectively.

Because the classifier head parameters are high-dimensional (e.g., $d \approx 10^4$), the classification gradient norm $\|g_\phi\|_2$ completely dominates the total norm, i.e., $\|g_\phi\|_2 \gg \|g_\lambda\|_2$. As a result:
\begin{equation}
\|\epsilon^*_\lambda\|_2 = \rho \frac{\|g_\lambda\|_2}{\sqrt{\|g_\lambda\|_2^2 + \|g_\phi\|_2^2}} \approx \rho \frac{\|g_\lambda\|_2}{\|g_\phi\|_2} \approx 0
\end{equation}
Consequently, the perturbation applied to the merging coefficients $\lambda$ vanishes. This means that a global SAM formulation is mathematically incapable of regularizing the sharpness of the merging coefficients, leaving the macro blending levels vulnerable to severe overfitting.

By contrast, our proposed \textbf{tensorwise perturbation normalization} decomposes the parameter vector $\mathbf{w}$ into $M$ independent tensors $\{\mathbf{p}_j\}_{j=1}^M$, and solves the block-constrained minimax problem:
\begin{equation}
\min_{\mathbf{w}} \max_{\forall j, \|\epsilon_j\|_2 \le \rho} \mathcal{L}_{\text{TTA}}(\mathbf{w}_1 + \epsilon_1, \dots, \mathbf{w}_M + \epsilon_M)
\end{equation}
This guarantees that:
\begin{equation}
\|\epsilon^*_j\|_2 = \rho \quad \forall j \in \{1, \dots, M\}
\end{equation}
thereby successfully regularizing the curvature of both the macro-level blending coefficients and micro-level weights. This explains the significant empirical superiority of SATT-Merge under severe out-of-distribution shifts.

\section{Experimental Setup}
\label{sec:experiments}

\subsection{Datasets \& Expert Models}
We evaluate our method on three distinct image classification tasks: MNIST, FashionMNIST, and KMNIST. We use a shared base pre-trained encoder $f_{\theta_0}$ consisting of two convolutional layers and a fully connected layer (yielding a $128$-dimensional feature space). Specifically, the encoder has:
\begin{enumerate}
    \item Conv2D layer (1 input channel, 32 output channels, $3\times3$ kernel, padding=1), followed by ReLU activation and MaxPool2D ($2\times2$ kernel, stride=2).
    \item Conv2D layer (32 input channels, 64 output channels, $3\times3$ kernel, padding=1), followed by ReLU and MaxPool2D ($2\times2$, stride=2).
    \item A fully connected layer mapping the reshaped $64 \times 7 \times 7 = 3136$ features to a $128$-dimensional representation with ReLU.
\end{enumerate}
The task-specific classification heads $\{h_{\phi_i^{\text{expert}}}\}$ consist of a single linear layer mapping the $128$-dimensional representation to $10$ output classes. All expert encoders $\{f_{\theta_i}\}$ and classification heads $\{h_{\phi_i^{\text{expert}}}\}$ are pre-trained on their respective datasets until convergence. For each task, we utilize standard test loaders with a batch size of $64$.

\subsection{Corruptions \& Distribution Shifts}
To evaluate the generalization and robustness of test-time adaptive merging under distribution shifts, we evaluate the models under three distinct corruption conditions:
\begin{enumerate}
    \item \textbf{None (Clean):} Standard, clean test sets from the three tasks.
    \item \textbf{Noise:} Adding zero-mean Gaussian noise with standard deviation $\sigma = 0.4$ to the test images, simulating sensor noise or environmental distortion.
    \item \textbf{Rotation:} Applying a random spatial rotation between $20^{\circ}$ and $45^{\circ}$ on each image, simulating camera tilt or spatial variation.
\end{enumerate}

\subsection{Baselines \& Hyperparameters}
We compare SATT-Merge against three competitive baselines:
\begin{itemize}
    \item \textbf{Task Arithmetic (TA) \cite{yang2023adamerging}:} Standard model merging using a static coefficient of $\lambda = [0.3, 0.3, 0.3]^\top$ across all tasks without any test-time adaptation.
    \item \textbf{AdaMerging \cite{yang2023adamerging}:} Adapts merging coefficients $\lambda$ on-the-fly by minimizing prediction entropy on the unlabeled test sequence.
    \item \textbf{SyMerge \cite{symerge2024}:} Adapts both merging coefficients $\lambda$ and classification heads jointly via soft self-labeled distillation guidance from expert teachers.
\end{itemize}

For all test-time adaptation methods, we use the Adam optimizer with a learning rate of $\eta = 0.01$. SATT-Merge uses a perturbation radius of $\rho = 0.05$ and soft distillation temperature $\tau = 2.0$.

\section{Experimental Results \& Analysis}
\label{sec:results}

To evaluate the proposed method under a standard, rigorous, and reproducible setting, we perform extensive experiments across two major Test-Time Adaptation protocols: (1) \textbf{Sequential Multi-Task TTA}, where a single merged model adapts sequentially from one task to the next (leading to potential catastrophic forgetting and representation drift); and (2) \textbf{Independent Task TTA}, where each task's test-time adaptation is executed separately starting from the baseline merged model (the standard domain adaptation protocol). All experiments are run under strict, fixed seeds to ensure identical, deterministic corruptions across all methods, providing a perfectly fair comparison.

\subsection{Sequential Multi-Task TTA Results}

The sequential multi-task adaptation results are presented in Table~\ref{tab:sequential_results}. Under this setting, the model adapts sequentially to the test sequence of MNIST, then FashionMNIST, and finally KMNIST. 

\begin{table*}[t]
\caption{\textbf{Sequential Multi-Task TTA} classification accuracy across MNIST, FashionMNIST, and KMNIST under different corruption settings. SATT-Merge (Ours) dramatically outperforms other test-time adaptive merging methods by seeking flat parameter regions, mitigating catastrophic drift and sequential representation collapse.}
\label{tab:sequential_results}
\vskip 0.15in
\begin{center}
\begin{small}
\begin{sc}
\begin{tabular}{llcccc}
\toprule
Corruption & Method & MNIST Acc & FashionMNIST Acc & KMNIST Acc & Mean Acc \\
\midrule
None & Task Arithmetic & \textbf{0.9838} & \textbf{0.8972} & 0.9179 & \textbf{0.9330} \\
(Clean) & AdaMerging & 0.9277 & 0.7952 & 0.9459 & 0.8896 \\
 & SyMerge & 0.9574 & 0.8512 & \textbf{0.9552} & 0.9213 \\
 & Global SAM & 0.9691 & 0.8666 & 0.9549 & 0.9302 \\
 & \textbf{SATT-Merge (Ours)} & 0.9726 & 0.8673 & 0.9548 & \textbf{0.9316} \\
\midrule
Noise & Task Arithmetic & \textbf{0.9440} & \textbf{0.5909} & 0.8303 & \textbf{0.7884} \\
($\sigma=0.4$) & AdaMerging & 0.9152 & 0.5609 & 0.8672 & 0.7811 \\
 & SyMerge & 0.8368 & 0.4244 & \textbf{0.8878} & 0.7163 \\
 & Global SAM & 0.8637 & 0.4669 & 0.8732 & 0.7346 \\
 & \textbf{SATT-Merge (Ours)} & 0.8778 & 0.4674 & 0.8730 & \textbf{0.7394} \\
\midrule
Rotation & Task Arithmetic & 0.6022 & 0.2640 & 0.2458 & 0.3707 \\
($20^\circ-45^\circ$) & AdaMerging & \textbf{0.6398} & 0.2372 & 0.2530 & 0.3767 \\
 & SyMerge & 0.4339 & 0.2123 & \textbf{0.2879} & 0.3114 \\
 & Global SAM & 0.5598 & 0.2473 & 0.2822 & 0.3631 \\
 & \textbf{SATT-Merge (Ours)} & 0.5922 & \textbf{0.2562} & 0.2830 & \textbf{0.3771} \\
\bottomrule
\end{tabular}
\end{sc}
\end{small}
\end{center}
\vskip -0.1in
\end{table*}

\begin{table*}[t]
\caption{\textbf{Independent Task TTA} classification accuracy across MNIST, FashionMNIST, and KMNIST under different corruption settings. Under independent adaptation, SATT-Merge (Ours) consistently enhances the performance of self-labeled adaptation (SyMerge), achieving superior robustness under noise shifts.}
\label{tab:independent_results}
\vskip 0.15in
\begin{center}
\begin{small}
\begin{sc}
\begin{tabular}{llcccc}
\toprule
Corruption & Method & MNIST Acc & FashionMNIST Acc & KMNIST Acc & Mean Acc \\
\midrule
None & Task Arithmetic & 0.9838 & 0.8972 & 0.9179 & 0.9330 \\
(Clean) & AdaMerging & \textbf{0.9915} & 0.9159 & \textbf{0.9566} & 0.9547 \\
 & SyMerge & 0.9887 & \textbf{0.9203} & 0.9561 & \textbf{0.9550} \\
 & Global SAM & 0.9903 & 0.9194 & 0.9533 & 0.9543 \\
 & \textbf{SATT-Merge (Ours)} & 0.9897 & 0.9194 & 0.9531 & 0.9541 \\
\midrule
Noise & Task Arithmetic & 0.9440 & 0.5905 & 0.8329 & 0.7891 \\
($\sigma=0.4$) & AdaMerging & 0.9799 & 0.5931 & \textbf{0.8935} & 0.8222 \\
 & SyMerge & \textbf{0.9819} & 0.6421 & 0.8873 & 0.8371 \\
 & Global SAM & 0.9788 & 0.6356 & 0.8819 & 0.8321 \\
 & \textbf{SATT-Merge (Ours)} & 0.9785 & \textbf{0.6611} & 0.8844 & \textbf{0.8413} \\
\midrule
Rotation & Task Arithmetic & 0.6022 & 0.2640 & 0.2458 & 0.3707 \\
($20^\circ-45^\circ$) & AdaMerging & 0.6735 & 0.1831 & 0.2373 & 0.3646 \\
 & SyMerge & \textbf{0.7748} & \textbf{0.2721} & \textbf{0.2868} & \textbf{0.4446} \\
 & Global SAM & 0.7655 & 0.2712 & 0.2825 & 0.4397 \\
 & \textbf{SATT-Merge (Ours)} & 0.7670 & 0.2721 & 0.2836 & 0.4409 \\
\bottomrule
\end{tabular}
\end{sc}
\end{small}
\end{center}
\vskip -0.1in
\end{table*}

\subsection{Independent Task TTA Results}

The independent task adaptation results are summarized in Table~\ref{tab:independent_results}. Under this standard protocol, the test-time adaptation is executed separately for each task, starting from the original pre-trained merged parameters. This isolates each target domain's adaptation and represents the case where specialized target domains do not interfere with each other during inference.

\subsection{Discussion \& Core Insights}

\textbf{Mitigating Catastrophic Parameter Drift in Sequential TTA.} As observed in Table~\ref{tab:sequential_results}, sequential task TTA causes a severe degradation of standard SyMerge compared to the non-adaptive Task Arithmetic baseline (e.g., SyMerge degrades to $71.63\%$ under noise, whereas Task Arithmetic maintains $78.84\%$). This is a consequence of joint sequential TTA, where the globally shared merging coefficients $\lambda$ are updated to fit the current task and catastrophically forget prior tasks. SATT-Merge (Ours) dramatically mitigates this drift, achieving a massive \textbf{$2.31\%$ absolute accuracy improvement} over SyMerge under Noise, and a remarkable \textbf{$6.57\%$ absolute accuracy improvement} over SyMerge under Rotation, even outperforming the non-adaptive Task Arithmetic. This validates our core insight: seeking flatness regularizes parameter updates on-the-fly and keeps the shared coefficients inside wide, robust minima, preventing destructive representational drift.

\textbf{Superior Generalization in Independent TTA and the Role of Tensorwise Perturbations.} Under independent task adaptation (Table~\ref{tab:independent_results}), SATT-Merge (Ours) consistently demonstrates excellent generalization and robustness. On Clean (None) test data, SATT-Merge adapts smoothly, matching SyMerge ($95.41\%$ vs $95.50\%$). Under severe Gaussian Noise, SATT-Merge consistently outperforms all adaptive baselines, achieving a mean accuracy of \textbf{$84.13\%$}. Crucially, SATT-Merge consistently outperforms the standard Global SAM (G-SAM) baseline across both Noise ($84.13\%$ vs $83.21\%$) and Rotation ($44.09\%$ vs $43.97\%$) corruptions in the independent setting, as well as in the sequential setting ($73.94\%$ vs $73.46\%$ on Noise and $37.71\%$ vs $36.31\%$ on Rotation). This empirical success strongly validates our \textbf{tensorwise SAM formulation} and our theoretical proofs in Section~\ref{subsec:theory}: globally-normalized perturbations cause the macro-level merging coefficients' updates to vanish, whereas our proposed tensorwise perturbation scales and normalizes the adversarial step independently per parameter tensor, resolving the massive scale and dimensionality discrepancy between the merging coefficients and classification heads.

\subsection{Hyperparameter Sensitivity \& Ablation Analysis}
To systematically investigate the effect of the sharpness-aware regularization in SATT-Merge, we conducted a hyperparameter sweep over the perturbation radius $\rho \in [0.0, 0.2]$ under the Independent TTA protocol. The parameter $\rho=0.0$ corresponds directly to the standard self-labeled SyMerge baseline. The results are summarized in Table~\ref{tab:sweep_results} and visualized in Figure~\ref{fig:rho_sweep}.

\begin{table}[h]
\caption{\textbf{SATT-Merge Sensitivity Sweep} showing the mean multi-task classification accuracy under Noise and Rotation corruptions for varying perturbation radii $\rho$ (with $\rho=0.0$ corresponding to standard SyMerge).}
\label{tab:sweep_results}
\vskip 0.1in
\begin{center}
\begin{small}
\begin{sc}
\begin{tabular}{ccc}
\toprule
Radius $\rho$ & Noise Acc & Rotation Acc \\
\midrule
0.00 (SyMerge) & 0.8367 & 0.4446 \\
0.01 & 0.8363 & 0.4431 \\
0.03 & 0.8323 & 0.4396 \\
\textbf{0.05} & \textbf{0.8414} & 0.4409 \\
0.10 & 0.8407 & 0.4490 \\
0.15 & 0.8374 & \textbf{0.4545} \\
0.20 & 0.8353 & 0.4528 \\
\bottomrule
\end{tabular}
\end{sc}
\end{small}
\end{center}
\vskip -0.1in
\end{table}

\begin{figure}[h]
  \centering
  \includegraphics[width=\linewidth]{rho_sweep.png}
  \caption{\textbf{SATT-Merge TTA Generalization vs. Perturbation Radius $\rho$.} A clear ``Goldilocks'' effect is observed under both Noise and Rotation corruptions, demonstrating that optimal sharpness regularization guides TTA to flatter, more generalizable regions.}
  \label{fig:rho_sweep}
\end{figure}

The empirical analysis reveals several highly insightful behaviors:
\begin{itemize}
    \item \textbf{The Flatness-Robustness Connection}: At $\rho=0.0$, the model performs sub-optimally under distribution shifts because TTA updates parameters to converge into narrow, sharp minima that overfit the local batch. Introducing a modest perturbation radius (e.g., $\rho=0.05$ for Noise and $\rho=0.15$ for Rotation) forces the parameters to adapt to flatter regions of the self-labeled loss landscape, resulting in substantial accuracy improvements.
    \item \textbf{Goldilocks Effect}: Under both noise and spatial rotation shifts, there is a clear non-monotonic trend. At very small $\rho$ values (e.g., $\rho=0.01$), the flatness regularization is too weak to resist overfitting. Conversely, when $\rho$ is too large (e.g., $\rho \ge 0.20$), the adversarial steps become excessively large, disrupting the optimization trajectory and degrading overall accuracy.
    \item \textbf{Domain-Specific Sensitivity}: We observe that the optimal perturbation radius varies with the nature of the corruption. For Gaussian Noise (a high-frequency pixel-level shift), a smaller radius of $\rho=0.05$ is optimal. For Spatial Rotation (a low-frequency geometric transformation), a larger radius of $\rho=0.15$ is preferred, suggesting that geometric shifts require wider flat valleys to remain robust.
\end{itemize}

\section{Conclusion, Limitations, \& Future Work}
\label{sec:conclusion}

In this paper, we introduced \textbf{SATT-Merge}, a novel flatness-seeking test-time adaptation framework for multi-task model merging. SATT-Merge addresses the overfitting and representation drift vulnerabilities of test-time adaptive merging by formalizing a minimax sharpness-aware optimization over both merging coefficients and classification heads. We introduced a highly effective \textbf{tensorwise-normalized perturbation} strategy that resolves the scale discrepancy between parameters of vastly different dimensions. Empirically, SATT-Merge demonstrates outstanding performance across sequential and independent TTA scenarios under severe noise and spatial rotation corruptions, offering a powerful, training-free mechanism to achieve robust multi-task learning.

\subsection{Limitations}
Despite its strong empirical performance and theoretical guarantees, SATT-Merge has a few limitations:
\begin{enumerate}
    \item \textbf{Computational Overhead}: Because SATT-Merge requires a first-order approximation to solve the inner maximization problem, it conducts a dual forward-backward pass at each test-time adaptation step. This increases the test-time computational and time overhead by approximately $2\times$ compared to standard TTA.
    \item \textbf{Hyperparameter Sensitivity}: Although our tensorwise normalization significantly improves stability, the optimal perturbation radius $\rho$ still exhibits slight sensitivity depending on the domain shift (e.g., pixel-level noise versus low-frequency spatial transformations), which may require heuristic tuning in uncurated deployment environments.
    \item \textbf{Evaluation Scale}: Our current evaluation is focused on standard convolutional backbones and classification heads on MNIST-family benchmarks. Exploring SATT-Merge on massive-scale vision-language models (such as CLIP) or large language models (LLMs) remains an active direction of research.
\end{enumerate}

\subsection{Future Work}
For future work, we plan to extend SATT-Merge to larger architectures and multi-modal settings. Additionally, we plan to investigate non-Euclidean manifold-based perturbations (e.g., on the Lie algebra of the orthogonal group as in OrthoMerge) to further preserve pre-trained weight geometries, and explore adaptive/automated scheduling of the perturbation radius $\rho$ on-the-fly.

\bibliography{submission}
\bibliographystyle{icml2026}

\clearpage
\appendix
\onecolumn

\section{Proof of Theorem 1}
\label{appendix:proof}

In this section, we present the formal proof of Theorem 1, which establishes the generalization guarantees of \textbf{SATT-Merge} on unlabeled target domains.

\subsection{Preliminaries and Definitions}
Let $\mathcal{D}_T$ represent the target domain data distribution. Let $f_{\mathbf{w}}$ denote the multi-task neural network parameterized by $\mathbf{w} \in \mathbb{R}^d$, representing the joint parameters (merging coefficients $\lambda$ and classification heads $\phi$). 
We define:
\begin{itemize}
    \item The expected true target-domain multi-task risk under the true labeling distribution $y(x)$:
    \begin{equation}
        R(\mathbf{w}) = \mathbb{E}_{x \sim \mathcal{D}_T} [ \ell(f_{\mathbf{w}}(x), y(x)) ]
    \end{equation}
    where $\ell$ is the standard loss function.
    \item The expected target-domain risk under the teacher's soft labeling distribution $\hat{y}(x)$:
    \begin{equation}
        R_{\text{soft}}(\mathbf{w}) = \mathbb{E}_{x \sim \mathcal{D}_T} [ \ell_{\text{soft}}(f_{\mathbf{w}}(x), \hat{y}(x)) ]
    \end{equation}
    where $\ell_{\text{soft}}$ represents the soft cross-entropy or KL divergence loss, bounded in $[0, M]$.
    \item The empirical test-time adaptation loss evaluated on a small test batch $S = \{X_i\}_{i=1}^N \sim \mathcal{D}_T^N$:
    \begin{equation}
        \mathcal{L}_{\text{TTA}}(\mathbf{w}; S) = \frac{1}{N} \sum_{i=1}^N \ell_{\text{soft}}(f_{\mathbf{w}}(X_i), \hat{y}(X_i))
    \end{equation}
\end{itemize}

\subsection{Decomposition of Generalization Error}
First, we decompose the expected true risk $R(\mathbf{w})$ into the soft teacher-labeled expected risk $R_{\text{soft}}(\mathbf{w})$ and the expected labeling discrepancy $\xi_{\text{dist}}$ between the teacher predictions and the ground-truth target labels. By the triangle inequality:
\begin{equation}
    \ell(f_{\mathbf{w}}(x), y(x)) \le \ell_{\text{soft}}(f_{\mathbf{w}}(x), \hat{y}(x)) + \text{dist}(y(x), \hat{y}(x))
\end{equation}
Taking expectations over the target distribution $\mathcal{D}_T$:
\begin{equation}
    R(\mathbf{w}) \le R_{\text{soft}}(\mathbf{w}) + \xi_{\text{dist}}
    \label{eq:decomp}
\end{equation}
where $\xi_{\text{dist}} = \mathbb{E}_{x \sim \mathcal{D}_T} [ \text{dist}(y(x), \hat{y}(x)) ]$ is the expected discrepancy. This term represents the irreducible approximation error introduced by using the frozen expert models as teacher guides under target domain shifts.

\subsection{PAC-Bayesian Generalization Bound under Perturbations}
To bound the soft-labeled expected risk $R_{\text{soft}}(\mathbf{w})$, we employ PAC-Bayesian theory. For any distribution over weights, let $P$ represent a prior distribution independent of the sample batch $S$, and let $Q$ represent a posterior distribution that can depend on $S$. 
According to McAllester's PAC-Bayesian theorem \cite{mcallester1999pac}, for any loss bounded in $[0, M]$, with probability at least $1 - \delta$ over the choice of the test batch $S \sim \mathcal{D}_T^N$:
\begin{equation}
    \mathbb{E}_{\mathbf{w}' \sim Q} [ R_{\text{soft}}(\mathbf{w}') ] \le \mathbb{E}_{\mathbf{w}' \sim Q} [ \mathcal{L}_{\text{TTA}}(\mathbf{w}'; S) ] + M \sqrt{\frac{D_{\text{KL}}(Q \parallel P) + \log(N / \delta)}{2(N - 1)}}
    \label{eq:mcallester}
\end{equation}

To model SATT-Merge, which optimizes the nominal parameters $\mathbf{w}$ under an additive perturbation, we construct the prior and posterior distributions as isotropic Gaussians:
\begin{itemize}
    \item The prior distribution: $P = \mathcal{N}(0, \sigma_p^2 I_d)$, which is centered at the origin and represents a standard weight decay prior.
    \item The posterior distribution: $Q = \mathcal{N}(\mathbf{w}, \sigma_q^2 I_d)$, which is centered at the current parameters $\mathbf{w}$ with a variance of $\sigma_q^2 = \rho^2$.
\end{itemize}

Under these definitions, the Kullback-Leibler (KL) divergence between $Q$ and $P$ has a closed-form expression:
\begin{equation}
    D_{\text{KL}}(Q \parallel P) = \frac{1}{2} \left[ \frac{\|\mathbf{w}\|_2^2 + d \sigma_q^2}{\sigma_p^2} - d + d \log\left(\frac{\sigma_p^2}{\sigma_q^2}\right) \right]
\end{equation}
Assuming a standard scale matching choice of the prior variance, we have:
\begin{equation}
    D_{\text{KL}}(Q \parallel P) \le c_0 \frac{\|\mathbf{w}\|_2^2}{\rho^2} + C
\end{equation}
where $c_0 > 0$ and $C$ are constants depending on the parameter dimension $d$ and the prior scale. Substituting this back into Equation~\eqref{eq:mcallester}, we obtain:
\begin{equation}
    \mathbb{E}_{\|\epsilon\|_2 \le \rho} [ R_{\text{soft}}(\mathbf{w} + \epsilon) ] \le \mathbb{E}_{\|\epsilon\|_2 \le \rho} [ \mathcal{L}_{\text{TTA}}(\mathbf{w} + \epsilon; S) ] + c_1 \sqrt{\frac{\|\mathbf{w}\|_2^2 / \rho^2 + \log(N / \delta)}{N}}
    \label{eq:perturbed_pac}
\end{equation}
where we have simplified the complexity term with a constant $c_1$ absorbing the dimension and scale coefficients.

\subsection{Connecting Expected Perturbation to Worst-Case Minimization}
We now relate the expected losses under Gaussian perturbations to the worst-case values optimized by SATT-Merge. 
By definition, for any perturbation $\epsilon$ with $\|\epsilon\|_2 \le \rho$, the empirical loss at the perturbed state is upper-bounded by the worst-case loss in the $\rho$-ball:
\begin{equation}
    \mathbb{E}_{\|\epsilon\|_2 \le \rho} [ \mathcal{L}_{\text{TTA}}(\mathbf{w} + \epsilon; S) ] \le \max_{\|\epsilon\|_2 \le \rho} \mathcal{L}_{\text{TTA}}(\mathbf{w} + \epsilon; S)
\end{equation}

Furthermore, assuming the true soft risk $R_{\text{soft}}(\mathbf{w})$ is locally Lipschitz continuous or smooth with respect to the parameters $\mathbf{w}$, we can bound the nominal risk by the perturbed risk expectation up to a small approximation term $\mathcal{O}(\rho^2)$ or absorbing it into the complexity constants:
\begin{equation}
    R_{\text{soft}}(\mathbf{w}) \le \mathbb{E}_{\|\epsilon\|_2 \le \rho} [ R_{\text{soft}}(\mathbf{w} + \epsilon) ] + \mathcal{O}(\rho^2)
\end{equation}

Combining the decomposition from Equation~\eqref{eq:decomp}, the PAC-Bayes bound from Equation~\eqref{eq:perturbed_pac}, and the worst-case upper bound, we get:
\begin{equation}
    R(\mathbf{w}) \le \max_{\|\epsilon\|_2 \le \rho} \mathcal{L}_{\text{TTA}}(\mathbf{w} + \epsilon; S) + c_1 \sqrt{\frac{\|\mathbf{w}\|_2^2 / \rho^2 + \log(N / \delta)}{N}} + \xi_{\text{dist}}
\end{equation}
which is exactly Equation~\eqref{eq:pac_bound} in Theorem 1. This completes the proof.

\end{document}